\subsection{Block Matrix Algebra}
Throughout this subsection, all vector spaces and matrices are over the same fixed field $\mathbb{F}$, and every matrix size is a positive integer.

\begin{definition} \label{def:block-partitions-and-conformality} \hlt{Block Partitions and Conformality}\\
Let $r,s\geq1$, let $m_1,\ldots,m_r$ and $n_1,\ldots,n_s$ be positive integers, and let $A\in\operatorname{Mat}\!\left(\sum\limits_{i=1}^{r}m_i,\sum\limits_{j=1}^{s}n_j,\mathbb{F}\right)$.\\
An \hlt{$r\times s$ block partition} $A=[A_{ij}]$ groups consecutive rows according to $(m_1,\ldots,m_r)$ and consecutive columns according to $(n_1,\ldots,n_s)$, so $(\forall\,i\in\{1,\ldots,r\})(\forall\,j\in\{1,\ldots,s\})\,A_{ij}\in\operatorname{Mat}(m_i,n_j,\mathbb{F})$. The submatrices $A_{ij}$ are the \hlt{blocks}.
\begin{enumerate}[label=(\roman*)]
\setlength{\itemsep}{0pt}
\item Two $r\times s$ block partitions $A=[A_{ij}]$ and $C=[C_{ij}]$ are \hlt{conformal for addition} if corresponding blocks have the same size.
\item An $r\times s$ block partition $A=[A_{ik}]$ and an $s\times t$ block partition $B=[B_{kj}]$, where $t\geq1$, are \hlt{conformal for multiplication} if\\
$\exists\,p_1,\ldots,p_t\in\mathbb{Z}_{>0}$ such that $A_{ik}\in\operatorname{Mat}(m_i,n_k,\mathbb{F})$ and $B_{kj}\in\operatorname{Mat}(n_k,p_j,\mathbb{F})$ $\forall i\in\{1,\ldots,r\}$, $k\in\{1,\ldots,s\}$, $j\in\{1,\ldots,t\}$.
\end{enumerate}
\end{definition}

\begin{definition} \label{def:adapted-block-matrix-of-linear-map} \hlt{Block Matrix of a Linear Map}\\
Suppose $\mathcal{V}=\mathcal{U}_1\oplus\mathcal{U}_2$ and $\mathcal{W}=\mathcal{Z}_1\oplus\mathcal{Z}_2$, where $\dim(\mathcal{U}_j)=n_j\geq1$ and $\dim(\mathcal{Z}_i)=m_i\geq1$ for $i,j\in\{1,2\}$. Choose ordered bases $\alpha_j$ of $\mathcal{U}_j$ and $\beta_i$ of $\mathcal{Z}_i$. By Definition \ref{def:adapted-basis-direct-sum}, $\alpha=(\alpha_1,\alpha_2)$ and $\beta=(\beta_1,\beta_2)$ are bases of $\mathcal{V}$ and $\mathcal{W}$ adapted to the decompositions.\\
For $T\in\LL(\mathcal{V},\mathcal{W})$, its representation matrix has conformal partition $[T]_{\beta\leftarrow\alpha}=\left[\begin{smallmatrix}A_{11}&A_{12}\\A_{21}&A_{22}\end{smallmatrix}\right]$, where $A_{ij}\in\operatorname{Mat}(m_i,n_j,\mathbb{F})$ for $i,j\in\{1,2\}$.\\
The $j$th block column contains coordinate columns of images of basis $\alpha_j$, while the $i$th block row retains their $\beta_i$-coordinate entries.
\end{definition}

\begin{proposition} \label{prop:block-matrix-arithmetic} \hlt{Block Matrix Arithmetic}\\
The following identities hold.
\begin{enumerate}[label=(\roman*)]
\setlength{\itemsep}{0pt}
\item If $A=[A_{ij}]$ and $C=[C_{ij}]$ are conformal $r\times s$ block partitions, then $A+C=[A_{ij}+C_{ij}]$.\\
In particular, conformal $2\times2$ partitions satisfy $A+C=\left[\begin{smallmatrix}A_{11}+C_{11}&A_{12}+C_{12}\\A_{21}+C_{21}&A_{22}+C_{22}\end{smallmatrix}\right]$.
\item If $A=[A_{ik}]$ and $B=[B_{kj}]$ are conformal for multiplication, then $(AB)_{ij}=\sum\limits_{k=1}^{s}A_{ik}B_{kj}$ for every admissible $i,j$.\\
In particular, conformal $2\times2$ partitions satisfy $AB=\left[\begin{smallmatrix}A_{11}B_{11}+A_{12}B_{21}&A_{11}B_{12}+A_{12}B_{22}\\A_{21}B_{11}+A_{22}B_{21}&A_{21}B_{12}+A_{22}B_{22}\end{smallmatrix}\right]$.
\end{enumerate}
\begin{proof}
\begin{enumerate}[label=(\roman*)]
\setlength{\itemsep}{0pt}
\item By Definition \ref{def:matrices-and-matrix-operations}, matrix addition is entrywise. Because conformality assigns the same entry positions to corresponding blocks, addition may be performed blockwise, so $A+C=[A_{ij}+C_{ij}]$. The case $r=s=2$ is the stated specialization.
\item By Proposition \ref{prop:matrix-product-entry-row-column}, each entry of $AB$ is the row-column sum over the shared index. Group that index into the $s$ consecutive intervals determined by the common block sizes. The contribution from the $k$th interval is the corresponding entry of $A_{ik}B_{kj}$.\\
Therefore $(AB)_{ij}=\sum\limits_{k=1}^{s}A_{ik}B_{kj}$. The case of two block rows and two block columns is the stated specialization.
\end{enumerate}
\hlt{Structural Commentary}: Addition is entrywise; multiplication groups the shared index.\\
\hlt{Key Insight}: Block arithmetic is ordinary matrix arithmetic performed at a coarser scale.
\end{proof}
\end{proposition}

\begin{proposition} \label{prop:stacked-nullspaces-and-concatenated-column-spaces} \hlt{Stacks Intersect Kernels, Concatenations Sum Columns}\\
Let $A\in\operatorname{Mat}(m,n,\mathbb{F})$, $B\in\operatorname{Mat}(p,n,\mathbb{F})$, $C\in\operatorname{Mat}(m,r,\mathbb{F})$, and $D\in\operatorname{Mat}(m,s,\mathbb{F})$. Then:
\begin{enumerate}[label=(\roman*)]
\setlength{\itemsep}{0pt}
\item $\ker\!\left(\left[\begin{smallmatrix}A\\B\end{smallmatrix}\right]\right)=\ker(A)\cap\ker(B)$.
\item $\operatorname{col}\!\left(\left[\begin{smallmatrix}C&D\end{smallmatrix}\right]\right)=\operatorname{col}(C)+\operatorname{col}(D)$.
\end{enumerate}
\begin{proof}
\begin{enumerate}[label=(\roman*)]
\setlength{\itemsep}{0pt}
\item By Proposition \ref{prop:block-matrix-arithmetic}, $\left[\begin{smallmatrix}A\\B\end{smallmatrix}\right]\mathbf{x}=\left[\begin{smallmatrix}A\mathbf{x}\\B\mathbf{x}\end{smallmatrix}\right]$ for every $\mathbf{x}\in\mathbb{F}^n$.\\
By Definition \ref{def:column-space-matrix-rank-nullity}, $\mathbf{x}\in\ker\!\left(\left[\begin{smallmatrix}A\\B\end{smallmatrix}\right]\right)$ if and only if $A\mathbf{x}=\mathbf{0}$ and $B\mathbf{x}=\mathbf{0}$, which is equivalent to $\mathbf{x}\in\ker(A)\cap\ker(B)$.
\item By Proposition \ref{prop:block-matrix-arithmetic}, $\left[\begin{smallmatrix}C&D\end{smallmatrix}\right]\left[\begin{smallmatrix}\mathbf{u}\\\mathbf{v}\end{smallmatrix}\right]=C\mathbf{u}+D\mathbf{v}$ for every $\mathbf{u}\in\mathbb{F}^r$ and $\mathbf{v}\in\mathbb{F}^s$.\\
By Definition \ref{def:column-space-matrix-rank-nullity}, $\operatorname{col}\!\left(\left[\begin{smallmatrix}C&D\end{smallmatrix}\right]\right)=\{C\mathbf{u}+D\mathbf{v}:\mathbf{u}\in\mathbb{F}^r,\ \mathbf{v}\in\mathbb{F}^s\}=\operatorname{col}(C)+\operatorname{col}(D)$.
\end{enumerate}
\hlt{Structural Commentary}: The kernel and column-space definitions translate both block actions.\\
\hlt{Key Insight}: Vertical blocks intersect kernels; horizontal blocks add column spaces.
\end{proof}
\end{proposition}

\begin{proposition} \label{prop:block-centrosymmetric-commutation} \hlt{Block Centrosymmetry and Commutation}\\
Let $n\geq1$, let $A,B,C,D\in\operatorname{Mat}(n,n,\mathbb{F})$, and set $M=\left[\begin{smallmatrix}A&B\\C&D\end{smallmatrix}\right]$ and $L_{2n}=\left[\begin{smallmatrix}0&I_n\\I_n&0\end{smallmatrix}\right]$ in $\operatorname{Mat}(2n,2n,\mathbb{F})$.\\
Then $M$ has \hlt{block centrosymmetric form}, meaning $A=D$ and $B=C$, if and only if $L_{2n}M=ML_{2n}$.
\begin{proof}
By Proposition \ref{prop:block-matrix-arithmetic}, $L_{2n}M=\left[\begin{smallmatrix}C&D\\A&B\end{smallmatrix}\right]$ and $ML_{2n}=\left[\begin{smallmatrix}B&A\\D&C\end{smallmatrix}\right]$.
\begin{enumerate}[label=\roman*.]
\setlength{\itemsep}{0pt}
\item[$(\Rightarrow)$] If $M$ has block centrosymmetric form, then $A=D$ and $B=C$, so the two block products agree. Therefore $L_{2n}M=ML_{2n}$.
\item[$(\Leftarrow)$] If $L_{2n}M=ML_{2n}$, then corresponding blocks are equal by Definition \ref{def:matrices-and-matrix-operations}.\\
Thus $C=B$ and $D=A$, so $M$ has block centrosymmetric form.
\end{enumerate}
\hlt{Structural Commentary}: Explicit block multiplication reduces commutation with $L_{2n}$ to equality of opposite blocks.\\
\hlt{Key Insight}: The exchange matrix detects block centrosymmetry.
\end{proof}
\end{proposition}

\begin{definition} \label{def:matrix-direct-sum} \hlt{Direct Sum of Matrices}\\
For $A\in\operatorname{Mat}(m,n,\mathbb{F})$ and $B\in\operatorname{Mat}(p,q,\mathbb{F})$, the \hlt{direct sum} of $A$ and $B$ is $A\oplus B\defeq\left[\begin{smallmatrix}A&0\\0&B\end{smallmatrix}\right]\in\operatorname{Mat}(m+p,n+q,\mathbb{F})$, where the upper-right and lower-left zero blocks have sizes $m\times q$ and $p\times n$, respectively.
\end{definition}
